% Topic 1.2.12 — Hypothesis Testing: Errors and Power.

\section{Hypothesis Testing: Errors and Power}
\label{sec:1.2.12-hypothesis-testing}

\begin{ticket}
Choice of statistical hypotheses. Simple and composite hypotheses. Statistical decisions and decision rules. Randomized decision rules. Type-I / type-II errors; power of a test. Most powerful and uniformly most powerful tests.
\end{ticket}

\subsection*{Theory}

A \emph{statistical hypothesis} is a claim about which probability
law on the sample space generated the observed data. Testing means
choosing, on the basis of data, between two competing hypotheses
while controlling the rates at which we make wrong decisions.

\medskip
\noindent\textbf{Statistical model and hypotheses.}

\begin{definition}[Hypothesis]
\label{def:hypothesis}
A \emph{statistical model} is a family of probability measures
$\{\Prob_\theta : \theta \in \Theta\}$ on a sample space
$(\mathcal{X}, \mathcal{A})$ indexed by a parameter
$\theta \in \Theta$. A pair of \emph{hypotheses}
$(H_0, H_1)$ is a partition $\Theta = \Theta_0 \sqcup \Theta_1$
into a \emph{null} hypothesis $H_0 : \theta \in \Theta_0$ and an
\emph{alternative} $H_1 : \theta \in \Theta_1$. A hypothesis is
\emph{simple} if its parameter set is a singleton; otherwise it is
\emph{composite}.
\end{definition}

\medskip
\noindent\textbf{Tests, errors, and power.}

\begin{definition}[Decision rule]
\label{def:decision-rule}
A \emph{(non-randomised) decision rule} or \emph{test} is a
measurable function $\varphi \colon \mathcal{X} \to \{0, 1\}$,
where $\varphi(x) = 1$ means ``reject $H_0$'' and $\varphi(x) = 0$
means ``do not reject''. The \emph{rejection region} (or
\emph{critical region}) is
$\mathcal{R} = \{x \in \mathcal{X} : \varphi(x) = 1\}$. A
\emph{randomised decision rule} replaces the codomain by $[0, 1]$:
$\varphi(x)$ is the conditional probability of rejecting given
$x$ is observed.
\end{definition}

\begin{definition}[Errors and power]
\label{def:errors-power}
For a (possibly randomised) test $\varphi$ in the model above, the
\emph{power function} is
\[
  \beta_\varphi(\theta) \;=\; \E_\theta\, \varphi \;=\; \Prob_\theta(\text{reject } H_0).
\]
\begin{itemize}
  \item For $\theta \in \Theta_0$, $\beta_\varphi(\theta)$ is the
        probability of \emph{type-I error} (incorrectly rejecting
        a true null).
  \item For $\theta \in \Theta_1$, the
        \emph{type-II error probability} is
        $1 - \beta_\varphi(\theta) = \Prob_\theta(\text{do not reject})$.
        The \emph{power} of $\varphi$ at $\theta$ is
        $\beta_\varphi(\theta)$ itself.
\end{itemize}
The \emph{size} (or \emph{level}) of the test is
$\sup_{\theta \in \Theta_0} \beta_\varphi(\theta)$.
\end{definition}

\medskip
\noindent\textbf{Most powerful and uniformly most powerful tests.}

\begin{definition}[(Uniformly) most powerful test]
\label{def:ump}
A test $\varphi^*$ of size~$\alpha \in (0, 1)$ is
\emph{most powerful (MP)} for testing $H_0 : \theta \in \Theta_0$
versus a simple alternative $H_1 : \theta = \theta_1$ if for every
test $\varphi$ of size at most~$\alpha$,
$\beta_\varphi(\theta_1) \leq \beta_{\varphi^*}(\theta_1)$. It is
\emph{uniformly most powerful (UMP)} of size~$\alpha$ for a composite
alternative $H_1 : \theta \in \Theta_1$ if the inequality above
holds for every $\theta_1 \in \Theta_1$ simultaneously.
\end{definition}

\begin{theorem}[Neyman--Pearson lemma]
\label{thm:neyman-pearson}
Consider testing the simple null $H_0 : \theta = \theta_0$ against
the simple alternative $H_1 : \theta = \theta_1$, with densities
$f_0$ and $f_1$ with respect to a common dominating measure on
$\mathcal{X}$. For each $k \geq 0$, the \emph{likelihood-ratio
test}
\[
  \varphi_k(x) \;=\; \begin{cases}
    1 & \text{if } f_1(x) > k\, f_0(x), \\
    0 & \text{if } f_1(x) < k\, f_0(x),
  \end{cases}
\]
(with $\varphi_k$ on the equality set $\{f_1 = k f_0\}$ chosen to
attain a desired size~$\alpha \in (0, 1)$) is most powerful of
size~$\alpha$ in the sense of \cref{def:ump}. Conversely, every
size-$\alpha$ MP test agrees with $\varphi_k$ for some $k \geq 0$
$\Prob_{\theta_0}$- and $\Prob_{\theta_1}$-a.s.\ outside the
equality set.
\end{theorem}
\proofof{neyman-pearson}

\begin{remark}
The lemma reduces optimal binary testing to a likelihood-ratio
threshold. When the alternative is composite, a UMP test exists in
nice cases (e.g., one-parameter exponential families with monotone
likelihood ratio) but not in general; for two-sided alternatives
in a Gaussian model the conventional ``equal-tail'' test is not UMP
and one falls back to other criteria — unbiasedness, invariance,
or likelihood-ratio against the entire alternative
\cite[Ch.~VII]{Shiryaev1989}.
\end{remark}

\medskip
\noindent\textbf{The standard $z$- and $t$-tests.} The Gaussian
sampling theorem of \cref{thm:sample-mean-variance} produces the
two most-used tests in classical statistics.

\begin{proposition}[Gaussian-mean tests]
\label{prop:z-and-t-tests}
Let $X_1, \dots, X_n$ be i.i.d.\ $\mathcal{N}(\mu, \sigma^2)$.
\begin{enumerate}[label=(\alph*)]
  \item \emph{Known $\sigma$.} For testing $H_0 : \mu = \mu_0$
        versus $H_1 : \mu > \mu_0$, the \emph{$z$-test} that
        rejects when $z = \sqrt n (\bar X - \mu_0)/\sigma > z_\alpha$
        (the upper-$\alpha$ quantile of $\mathcal{N}(0, 1)$) is UMP
        of size~$\alpha$.
  \item \emph{Unknown $\sigma$.} For the same hypotheses, the
        \emph{$t$-test} that rejects when
        $t = \sqrt n (\bar X - \mu_0)/S > t_{\alpha, n - 1}$
        (the upper-$\alpha$ quantile of $t_{n - 1}$) has size
        exactly~$\alpha$ under $H_0$, with the test statistic
        following $t_{n - 1}$ by \cref{thm:sample-mean-variance}.
\end{enumerate}
The two-sided versions reject when $|z| > z_{\alpha/2}$ or
$|t| > t_{\alpha/2, n - 1}$.
\end{proposition}
\proofof{z-and-t-tests}

\subsection*{Proofs}

\begin{proof}[\Cref{thm:neyman-pearson}, sketch]
\proofanchor{neyman-pearson}
Suppose $\varphi^* = \varphi_k$ has size~$\alpha$ — i.e.,
$\E_{\theta_0} \varphi^* = \alpha$ — and let $\varphi$ be any test
of size $\leq \alpha$. Consider the integral of
$(\varphi^* - \varphi)(f_1 - k f_0)$ over $\mathcal{X}$. By the
definition of $\varphi^*$, the integrand is non-negative pointwise:
\begin{itemize}
  \item on $\{f_1 > k f_0\}$, $\varphi^* = 1 \geq \varphi$ and $f_1 - k f_0 > 0$;
  \item on $\{f_1 < k f_0\}$, $\varphi^* = 0 \leq \varphi$ and $f_1 - k f_0 < 0$;
  \item on $\{f_1 = k f_0\}$, $f_1 - k f_0 = 0$.
\end{itemize}
Integrating,
\[
  0 \leq \int (\varphi^* - \varphi)(f_1 - k f_0) = \bigl(\E_{\theta_1} \varphi^* - \E_{\theta_1} \varphi\bigr) - k \bigl(\E_{\theta_0} \varphi^* - \E_{\theta_0} \varphi\bigr).
\]
Since $\E_{\theta_0} \varphi^* = \alpha$ and $\E_{\theta_0} \varphi \leq \alpha$,
the second bracket is $\leq 0$, so $-k$ times it is $\geq 0$, and the
inequality $\geq 0$ on the whole expression then yields
$\E_{\theta_1} \varphi^* \geq \E_{\theta_1} \varphi$ — i.e., the
power at $\theta_1$ of $\varphi^*$ dominates that of any other test
of size $\leq \alpha$. The converse identifies the structure of any
optimal test: equality $\varphi = \varphi^*$ outside the
$\{f_1 = k f_0\}$ set follows from the strict-inequality cases above
because the integrand is strictly positive there and the integral
must be zero for a competing optimum.
\end{proof}

\begin{proof}[\Cref{prop:z-and-t-tests}]
\proofanchor{z-and-t-tests}
\emph{(a)} For known $\sigma$, the joint density of an
$\mathcal{N}(\mu, \sigma^2)$ sample factorises as
$f_\mu(\vect{x}) \propto \exp\bigl(\mu \sum_i x_i / \sigma^2 - n \mu^2 / (2 \sigma^2)\bigr)$
times a $\mu$-independent factor. The likelihood ratio
$f_{\mu_1}(\vect{x})/f_{\mu_0}(\vect{x})$ is therefore an increasing
function of $\sum_i x_i$ (equivalently, of $\bar X$) whenever
$\mu_1 > \mu_0$. By the Neyman--Pearson lemma
(\cref{thm:neyman-pearson}), for any specific $\mu_1 > \mu_0$ the
MP test rejects for large~$\bar X$; the threshold is the same
function of the size~$\alpha$ regardless of the specific $\mu_1$
chosen. So the test ``reject when $\bar X > c_\alpha$'' is MP for
every $\mu_1 > \mu_0$ simultaneously, hence UMP for the composite
alternative $\mu > \mu_0$. Standardising, $c_\alpha$ corresponds
to $z = \sqrt n (\bar X - \mu_0)/\sigma > z_\alpha$.

\emph{(b)} The size claim follows directly from
\cref{thm:sample-mean-variance}: under $H_0 : \mu = \mu_0$,
$T = \sqrt n (\bar X - \mu_0)/S \sim t_{n - 1}$, so
$\Prob_{\mu_0}(T > t_{\alpha, n - 1}) = \alpha$. (UMP-ness fails
for the $t$-test because the alternative $\mu > \mu_0$ now has the
nuisance parameter~$\sigma$; the standard $t$-test is UMP among
\emph{unbiased} or \emph{invariant} tests, by deeper arguments
\cite[Ch.~VII]{Shiryaev1989}.)
\end{proof}

\subsection*{Examples}

\begin{example}[$z$-test power calculation]
\label{ex:z-test-power}
Test $H_0 : \mu = 0$ against $H_1 : \mu > 0$ from $n = 25$
observations of an $\mathcal{N}(\mu, 1)$ sample at level
$\alpha = 0.05$. The UMP test (\cref{prop:z-and-t-tests}\,(a))
rejects when $\bar X > z_{0.05}/\sqrt{25} = 1.645/5 = 0.329$. The
power at $\mu = 0.5$ is
\[
  \beta_\varphi(0.5) = \Prob_{0.5}(\bar X > 0.329) = \Prob\bigl(Z > (0.329 - 0.5)\sqrt{25}\bigr) = \Prob(Z > -0.855) \approx 0.804,
\]
where $Z \sim \mathcal{N}(0, 1)$. So with $n = 25$ we detect a true
mean of $0.5$ with probability about~$0.80$.
\end{example}

\begin{example}[Sample-size determination]
\label{ex:sample-size}
Continuing the previous setup, what $n$ is needed to detect
$\mu = 0.5$ with power~$0.95$? The condition
\[
  1 - \Phi\bigl(z_\alpha - \mu_1 \sqrt n / \sigma\bigr) = 0.95
\]
becomes $z_\alpha - 0.5 \sqrt n = -z_{0.05} = -1.645$, i.e.\
$\sqrt n = (z_\alpha + 1.645)/0.5 = (1.645 + 1.645)/0.5 = 6.58$,
giving $n \geq 44$. The sample-size formula
$n = ((z_\alpha + z_\beta)\sigma/(\mu_1 - \mu_0))^2$ generalises this.
\end{example}

\begin{problem}
\label{prob:bernoulli-test}
$X_1, \dots, X_n$ are i.i.d.\ $\mathrm{Bernoulli}(p)$. Show that
the most powerful test of $H_0 : p = 1/2$ against $H_1 : p = 3/4$
at level $\alpha$ rejects when the number of successes $S = \sum X_i$
exceeds a threshold, and identify it.
\end{problem}

\begin{proof}[Solution]
The likelihood is $p^S (1 - p)^{n - S}$, so the likelihood ratio
$L(\vect{x}) = (3/4)^S (1/4)^{n - S} / (1/2)^n = 2^n \cdot 3^S / 4^n = 3^S / 2^n$ —
strictly increasing in~$S$. The Neyman--Pearson lemma
(\cref{thm:neyman-pearson}) therefore gives the MP test: reject
when $S > c$ for some integer threshold~$c$. Choose $c$ so that
$\Prob_{p = 1/2}(S > c) = \alpha$, i.e., $c$ is the upper
$\alpha$-critical value of $\mathrm{Bin}(n, 1/2)$. Concretely, for
$n = 20$ and $\alpha = 0.05$, the binomial CDF gives $c = 14$
($\Prob(S \geq 15 \mid p = 1/2) \approx 0.021$, while
$\Prob(S \geq 14 \mid p = 1/2) \approx 0.058$); a randomised test
on $\{S = 14\}$ would attain exactly $\alpha = 0.05$ and is
strictly MP. By the same monotone-likelihood-ratio reasoning, this
test is UMP for the one-sided composite alternative $p > 1/2$.
\end{proof}
